
\slajdid{05-s052}
\begin{frame}[label=05-s052]
\gusttekst\setlength{\abovedisplayskip}{3pt}\setlength{\belowdisplayskip}{3pt}\setlength{\parskip}{3pt}
\begin{center}ZAOKRUŽIVANJE REZULTATA MNOŽENJA\end{center}
Da bi se smanjila greška odsecanja u mnogim sistemima se radi zaokruživanje rezultata.
\[a_{n-1}.a_{n-2}\ldots a_1a_0\times b_{n-1}.b_{n-2}\ldots b_1b_0=\textcolor{DEcrvena}{c_{2n-1}.c_{2n-2}\ldots c_n}\textcolor{DEplava}{c_{n-1}c_{n-2}\ldots c_1c_0}\]
Vrednost $\textcolor{DEcrvena}{c_{2n-1}.c_{2n-2}\ldots c_n+r}$ se smešta u $n$ bitni registar i koriste za dalju obradu\\
$r$ je dodatna vrednost zaokruživanja i zavisi od načina koji se primenjuje.\\
Sigurno je manja ili jednaka od vrednosti težine najnižeg bita $\textcolor{DEcrvena}{c_n\equiv c_{-n+1}}$

Zaokruživanje na dole – prema minus beskonačno – praktično odsecanje

Zaokruživanje na gore – prema plus beskonačno – $r=\dfrac{1}{2^{n-1}}$\\
dodaje se vrednost bita najmanje težine

Zaokruživanje prema nuli

Zaokruživanje od nule

Zaokruživanje najbližoj vrednosti $\textcolor{DEplava}{r=c_{n-1}}$ pa onda prema parnoj, neparnoj itd…
\end{frame}
